\section*{Bab \ref{ch:g5:human-reproduction} --- Perkembangbiakan Manusia}

\begin{solution}{exo:g5:human-reproduction:1}
\omterm{def:g4:animal-reproduction:sexual}{Sel telur} dari ibunya, \omterm{def:g4:animal-reproduction:sexual}{sel sperma} dari ayahnya; penggabungan keduanya
adalah \omterm{def:g5:flower-to-fruit:steps}{pembuahan}.
\end{solution}

\begin{solution}{exo:g5:human-reproduction:2}
Di dalam \omterm{prop:g5:human-reproduction:plan}{rahim} ibunya, selama kira-kira sembilan bulan --- masa
\omterm{def:g5:human-reproduction:pregnancy}{kehamilannya}.
\end{solution}

\begin{solution}{exo:g5:human-reproduction:3}
\omterm{def:g5:human-reproduction:pregnancy}{Embrio} lebih dahulu, lalu \omterm{def:g5:human-reproduction:pregnancy}{janin} --- mulai kira-kira dua bulan, begitu
rancangan tubuhnya terbentang.
\end{solution}

\begin{solution}{exo:g5:human-reproduction:4}
Lewat \omterm{prop:g5:human-reproduction:cord}{tali pusar}: \omterm{def:g4:journey-of-food:digestion}{zat gizi} dan oksigen menyeberang dari \omterm{prop:g4:journey-of-food:absorb}{darah} ibunya
ke dalam \omterm{prop:g4:journey-of-food:absorb}{darah} bayinya, dan buangan bayinya menyeberang kembali untuk
ditangani tubuh ibunya.
\end{solution}

\begin{solution}{exo:g5:human-reproduction:5}
Tempat tali pusarmu sendiri dahulu melekat, sebelum ia dipotong pada
saat \omterm{def:g2:born-grow-die:cycle}{kelahiran}.
\end{solution}

\begin{solution}{exo:g5:human-reproduction:6}
Tangis pertamanya mengisi \omterm{def:g4:breathing:organs}{paru-parunya} dengan udara untuk pertama
kalinya: sampai saat itu \omterm{def:g4:breathing:organs}{paru-parunya} belum pernah bernapas ---
oksigennya selalu tiba lewat \omterm{prop:g5:human-reproduction:cord}{tali pusarnya}.
\end{solution}

\begin{solution}{exo:g5:human-reproduction:7}
\omterm{def:g5:flower-to-fruit:steps}{Pembuahan} di dalam tubuh ibunya; \omterm{def:g2:born-grow-die:cycle}{kelahiran} yang \omterm{def:g4:animal-reproduction:ovivivi}{vivipar} (bayinya
tumbuh di dalam \omterm{prop:g5:human-reproduction:plan}{rahim} lalu \omterm{def:g2:animals-grow-up:born}{lahir hidup}); pemberian \omterm{def:g2:animals-grow-up:born}{susu} sesudahnya ---
persis ciri kelas \omterm{prop:g3:sorting-living-things:groups}{mamalia}.
\end{solution}

\begin{solution}{exo:g5:human-reproduction:8}
Karena apa pun yang ada di dalam \omterm{prop:g4:journey-of-food:absorb}{darah} ibunya mencapai bayinya lewat
\omterm{prop:g5:human-reproduction:cord}{tali pusar} --- termasuk zat berbahaya dari asap dan alkohol --- dan
seorang \omterm{def:g5:human-reproduction:pregnancy}{janin} jauh lebih rapuh daripada orang dewasa.
\end{solution}

\begin{solution}{exo:g5:human-reproduction:9}
Kembar identik: satu \omterm{def:g4:animal-reproduction:sexual}{sel telur} terbuahi yang terbelah menjadi dua
\omterm{def:g5:human-reproduction:pregnancy}{embrio}. Kembar biasa (fraternal): dua \omterm{def:g4:animal-reproduction:sexual}{sel telur} yang dibuahi oleh dua
\omterm{def:g4:animal-reproduction:sexual}{sel sperma}, dan hanya berbagi \omterm{prop:g5:human-reproduction:plan}{rahim}. Pasangan yang identik serupa
bagai salinan karena sebuah tubuh dibangun dari \omterm{def:g5:first-look-at-cells:cell}{sel} pertamanya --- dan
keduanya bermula dari \omterm{def:g5:first-look-at-cells:cell}{sel} yang sama.
\end{solution}

\begin{solution}{exo:g5:human-reproduction:10}
Di ujung sang gajah, bahkan melampauinya: biasanya satu bayi,
bertahun-tahun ketidakberdayaan, dan perawatan yang memuat bukan hanya
pemberian makan dan perlindungan melainkan juga pengajaran --- bahasa,
keterampilan, sekolah. Jutaan-tanpa-perawatan milik \omterm{prop:g3:sorting-living-things:groups}{ikan} kod adalah
kutub yang berlawanan.
\end{solution}

\begin{solution}{exo:g5:human-reproduction:11}
Tugas memuat oksigen dan membongkar karbon dioksida: sampai
\omterm{def:g2:born-grow-die:cycle}{kelahirannya} \omterm{prop:g4:journey-of-food:absorb}{darah} ibunya yang mengerjakannya lewat \omterm{prop:g5:human-reproduction:cord}{tali pusar}; sejak
napas pertamanya, kantong udara \omterm{def:g4:breathing:organs}{paru-paru} bayinya sendiri yang
mengambil alih. Tangisnya tepat waktu karena ia datang persis pada
menit jalur pasokan yang lama dipotong --- \omterm{def:g4:breathing:organs}{paru-parunya} harus mulai
tepat ketika \omterm{prop:g5:human-reproduction:cord}{tali pusarnya} berhenti.
\end{solution}
